\documentclass{article}
\usepackage[T1]{fontenc}
\usepackage{lmodern}
\usepackage{graphicx}
\usepackage{amsmath,amssymb}
\usepackage[a4paper,margin=2cm]{geometry}
\usepackage[absolute,overlay]{textpos}
\setlength{\TPHorizModule}{1mm}
\setlength{\TPVertModule}{1mm}
\usepackage[indonesian]{babel}
\usepackage{hyperref}
\setlength{\emergencystretch}{2em}
% unit: Halaman 46

\begin{document}

% === Halaman 46 (llm) ===

\section*{DES Berganda}

\begin{itemize}
    \item Karena DES mempunyai potensi kelemahan pada \textit{brute force attack}, maka dibuat varian dari DES.
    \item Varian DES yang paling luas digunakan adalah DES berganda (\textit{multiple DES}).
    \item DES berganda adalah enkripsi berkali-kali dengan DES dan menggunakan kunci ganda.
    \item Tinjau dua macam DES berganda:
    \begin{enumerate}
        \item \textit{Double DES}
        \item \textit{Triple DES (3DES)}
    \end{enumerate}
\end{itemize}

\end{document}
